\section{Lean formalization: definitions, checked kernel, and deep targets}
\label{sec:formalization}

The formalization path has a useful small completed component and a much
larger explicitly specified endpoint. The completed component checks
finite-table locality and compilation. The endpoint is a family of
object-language \texttt{PRAAx} derivations. We do not identify the two.

\subsection{The checked finite-table kernel}

The file \texttt{lean/FiniteObservation.lean} imports \texttt{Std}, not
Mathlib, and was checked with Lean 4.32.0. Its central definitions are:
\begin{lstlisting}
abbrev Stream := Nat -> Nat
abbrev Table := List Nat

def PrefixEq (h : Nat) (a b : Stream) : Prop :=
  forall i, i <= h -> a i = b i

def tabulate (a : Stream) (h : Nat) : Table :=
  (List.range (h + 1)).map a

def demand (u : Nat -> Nat) : Nat -> Nat -> Nat
  | 0, k => k
  | n + 1, k => u (demand u n k)

def Local (u : Nat -> Nat)
    (t : Nat -> Stream -> Stream) : Prop :=
  forall s k a b,
    PrefixEq (u k) a b -> PrefixEq k (t s a) (t s b)

def ofBody (u : Nat -> Nat)
    (body : Nat -> Nat -> Table -> Nat)
    (s : Nat) (a : Stream) (p : Nat) : Nat :=
  body s p (tabulate a (u p))
\end{lstlisting}

The theorem \texttt{ofBody\_local} proves that a monotone numerical query
bound makes this table body local. The numerical budget at stage $s$ is
\[
 \mathtt{horizon}(u,N,k,s)=\mathtt{demand}(u,N-s,k).
\]
The executable \texttt{build} fixes $N,k$ and recurses on the stage,
storing a \texttt{List Nat}. The semantic \texttt{iter} recurses into
\texttt{Stream}; it is explicitly a reference specification, not a
claimed first-order implementation. The main checked theorem is
\begin{lstlisting}
theorem compiled_correct
    (ht : Local u t) (a : Stream) (N k : Nat) :
  lookup (build u t a N k N) k = iter t a N k
\end{lstlisting}
with the implicit parameters and types as in the file.

The inductive proof uses a stronger finite-prefix assertion at every
stage. At a successor, the budget identity
$h_s=u(h_{s+1})$ supplies exactly the input agreement needed by locality.
This is a real verification of the generic compilation mechanism, not a
theorem whose premise is the same as its conclusion.

The same file proves \texttt{allBelow\_iff}, the exact interpretation of
a finite Boolean conjunction; \texttt{finite\_safety}, a finite invariant
lemma; and \texttt{descent\_fuel}, a scalar strict-descent bound. The
latter is auxiliary infrastructure rather than a new Mordell--Weil
result. It states that a natural rank which decreases at every step for
$H$ steps satisfies $H\le\operatorname{rank}(0)$.

The axiom printouts for the four principal checked theorems contain
\texttt{propext} and \texttt{Quot.sound}. They contain neither
\texttt{Classical.choice} nor \texttt{sorryAx}. These are Lean's ambient
foundational dependencies; they are not axioms added to \iPRA. An axiom
audit alone would not establish the object-theory claim in either case.

\subsection{Type-checked mathematical specifications}

The separate file \texttt{lean/Targets.lean} defines the numerical and
rational interfaces for the main analytic and algebraic targets. Its
structures and predicates are specifications of proof obligations;
constructing a type or a structure with unfilled obligations is not a
proof that it is inhabited. No opaque theorem, new axiom, or
\texttt{sorry} is used to assert those conclusions. The file's comments
mark that status explicitly.

At this semantic specification layer it is convenient to use
\texttt{Rat}, lists, and Lean functions. Their object-language
counterparts must be natural codes and the fixed \texttt{PR} terms for
the algorithms in the manuscript. An arbitrary Lean function
\texttt{Nat -> Nat} carries no automatic evidence of primitive
recursiveness. That evidence is a separate deliverable at the next layer.

\subsection{The existing deep backend}

The inspected source files provide the following exact interfaces:
\begin{longtable}{@{}>{\raggedright\arraybackslash}p{0.35\linewidth}>{\raggedright\arraybackslash}p{0.59\linewidth}@{}}
\toprule
Object & Meaning\\
\midrule
\endhead
\texttt{PR k} & Syntax for $k$-ary primitive recursive functions,
with constructors \texttt{zero}, \texttt{succ}, \texttt{proj},
\texttt{comp}, \texttt{prec}.\\
\texttt{PTm n} & Terms in $n$ free numerical variables, with
\texttt{var} and \texttt{app}.\\
\texttt{Fm PTm n} & First-order formulas over these terms; formulas
are not Lean predicates in this layer.\\
\texttt{Prf PRAAx Gamma phi} & Intuitionistic derivability of the
encoded formula from the encoded context.\\
\texttt{PRA.qfInd} & Constructs an induction proof from an encoded
formula, a syntactic \texttt{QF} witness, and object-level base/step proofs.\\
\texttt{Prf.subst}, \texttt{Prf.weaken}, \texttt{Prf.trans} &
Capture-avoiding substitution, weakening, and simultaneous cut.\\
\texttt{pra\_iff\_native} & Transfers completed arithmetic proofs to
the native scalar calculus.\\
\bottomrule
\end{longtable}
These are read from the pinned sources \cite{repo-pr,repo-deduction,repo-native};
we have not rebuilt the full repository dependency closure in this run.

\subsection{The object-language plan, with no semantic shortcut}

\paragraph{Stage 1: numerical representations and PR definitions.}
Implement the pair, list, integer, and rational codes from
Section~\ref{sec:foundations}, including total default branches. For each
operation construct its actual syntax in \texttt{PR k}. A semantic
function interpreting these syntax trees can be defined in Lean to state
correctness, but it is a metatheoretic function. Do not postulate an
object-level PR evaluator for arbitrary \texttt{PR} syntax.

Define addition, truncated subtraction, products, powers, list lookup,
bounded folds, rational order, finite determinants, and polynomial
convolution through \texttt{comp} and \texttt{prec}. This is repetitive
syntax construction and can be generated, but the output terms and their
arithmetic defining-equation proofs must remain inspectable. The
universal object-level facts come from those proofs, not from running a
finite set of examples.

\paragraph{Stage 2: Boolean reflection as an object theorem.}
For a chosen finite predicate $B$, build a specific Boolean PR term
$\chi_B$. Encode $B\leftrightarrow\chi_B=1$ in
\texttt{Fm PTm n} and prove it using the fixed-formula construction of
Lemma~\ref{lem:reflection}. The careful two-direction proof of bounded
universal reflection avoids circularly using bounded induction before
establishing it. A syntactic QF checker certifies the formulas on which
\texttt{PRA.qfInd} is called.

\paragraph{Stage 3: finite safety and contracts.}
For each selected numerical invariant, form the exact formula
\[
 j\le H(\mathbf x)\longrightarrow
   I(\mathbf x,j,s_j(\mathbf x))=1.
\]
Its order test is a Boolean PR term, so the encoded formula is
quantifier-free. Supply the object derivations of the base and step,
then apply \texttt{PRA.qfInd}. Global name or continuity assumptions
may remain in the context. No rule infers the base or step merely from
an ambient Lean proof about standard natural numbers.

The generic compiler's local verification obligation is a single
arithmetic theorem with the stage, horizon, and metadata as free
numerical variables. This distinguishes a uniform theorem from a
metatheoretic family of proofs for separately fixed stages. Metadata
bounds must be monotone or replaced by bounded maxima, as in the
construction of $U_N$.

\paragraph{Stage 4: the table compiler.}
Construct PR syntax for the demand list, the initial observation table,
and the forward stage fold. Prove prefix locality for its fixed body,
then mirror the finite proof of \texttt{build\_correct}. At the deep
level its invariant is a Boolean conjunction over the finite rows,
not the semantic proposition about \texttt{Stream}s. Include the
initial-contract derivation. Include Cauchy regularity among the finite
checks whenever the output is supposed to be a regular name.

The resulting generic proof is a compiler \emph{for supplied arithmetic
proof packets}: numerical function definitions, formula definitions,
QF witnesses, and uniform local proof terms. It is not an unrestricted
reflection principle for arbitrary Lean mathematics.

\paragraph{Stage 5: factorization.}
The mathematical statements to encode are
Lemmas~\ref{lem:fta-descent-contract}, \ref{lem:fta-root-table}, and
\ref{lem:fta-deflation-table}. The imported descent algorithm requires
its own deep rational-arithmetic proof; a Lean section variable stating
its correctness is an explicit assumption until that work is done.
The additional tasks are dyadic rounding, the monic trajectory
invariant, normalized root regularity, deflation regularity, and the
four-term polynomial defect estimate. All function symbols and bounds
are displayed in Section~\ref{sec:fta-extension}.

After these local proofs, encode the cap and backward-horizon
recursions and apply the finite-table compiler. Prove the finite
coefficient telescoping identity and its $dWG^d2^{-t}$ bound. Fix $k$
before the stage induction and introduce its universal quantifier only
at the end. The final formula is exactly the rational coefficient
inequality \eqref{eq:fta-uniform-factorization}, together with the
regularity assertions, rather than an invocation of Lean's
\texttt{IsAlgClosed} for the ordinary complex field.

\paragraph{Stage 6: convexity and finite nets.}
Define dyadic grids as finite tuples, rational rounding, first-index
minimization, interval intersection, and the explicit schedules
$\mu(k+2)$, $v(k)$, and $2k+c$. Prove the finite sampled minimum
inequality first. Instantiate the continuity premise only at the rounded
point and the required finite precisions. The two endpoint secant
inequalities yield the Bohr--Mollerup modulus once finite AM--GM and
positive-power comparison have their arithmetic proofs.

The positivity witness is part of the original presentation. Rational
non-strict inequalities between names are verified by explicit finite
error bounds and Lemma~\ref{lem:vanishing-error}, not by assuming real
trichotomy. The canonical extension is constructed through the derived
modulus; separately assigned irrational values remain outside the claim.

\paragraph{Stage 7: finite algebraic sections.}
Define the rational polynomial and point-list predicates, synthetic
division, determinant enumeration, maximal-minor rank, and the explicit
annihilator. The finite polynomial root bound is an induction over a
supplied distinct root list. Then split on the decidable last coordinate
of the annihilator and prove the two branches of
Theorem~\ref{thm:superelliptic-section}. Implement subset-mask cover
verification and first-cover assignment. The cover theorem becomes an
arithmetic implication from a successful finite test. Its success on
all intended point tuples is a separate future mathematical theorem.

\paragraph{Stage 8: integration and proof audit.}
Only after the \texttt{Prf PRAAx} terms exist transfer them through
\texttt{pra\_iff\_native}. Record axiom printouts of the complete
endpoints and a build of their transitive dependencies. An integration
that leaves the desired conclusion in a structure field, an assumption,
an \texttt{opaque} declaration, or a \texttt{sorry} has not completed
this stage. The current archive explicitly stops before claiming it.

\subsection{Concrete proof targets and completion criteria}

\begin{longtable}{@{}>{\raggedright\arraybackslash}p{0.29\linewidth}>{\raggedright\arraybackslash}p{0.40\linewidth}>{\raggedright\arraybackslash}p{0.25\linewidth}@{}}
\toprule
Target & Required mathematical invariant & Current status\\
\midrule
\endhead
Prefix locality and table compilation & Equality on the scheduled
finite input prefixes implies equality of the output rows. & Ambient
Lean proof checked; deep coding pending.\\[4pt]
Finite Boolean reflection & Fold value one iff all bounded tests
succeed, with the noncircular induction proof. & Ambient finite-fold
lemma checked; general deep reflection pending.\\[4pt]
Uniform factorization & Good finite tables, degree/cap facts, local
coefficient defects, finite telescope. & Human proof from stated
Descent dependency; full Lean proof pending.\\[4pt]
Candidate continuity modulus & Finite AM--GM, witnessed positivity,
two rational secant inequalities. & Human proof; typed targets and
concrete arithmetic roadmap.\\[4pt]
Infimum and coherent optimizer & Finite minimum and cross-interval
checks; quantitative uniqueness at fixed precisions. & Human proof;
exact examples and typed targets.\\[4pt]
Superelliptic section bound & Quotient degree and remaining listed
roots; annihilator identity; finite grouping. & Human proof;
exact examples and typed targets.\\
\bottomrule
\end{longtable}
